\documentclass[UTF8,12pt,a4paper]{ctexart}

\usepackage[margin=25mm]{geometry}
\usepackage{booktabs}
\usepackage{enumitem}
\usepackage{hyperref}
\usepackage{listings}
\usepackage{tabularx}

\hypersetup{
  colorlinks=true,
  linkcolor=blue,
  urlcolor=blue,
  pdftitle={Chemex 使用手册},
  pdfauthor={Chemex},
}

\lstdefinestyle{chemexcode}{
  basicstyle=\ttfamily\small,
  columns=fullflexible,
  keepspaces=true,
  breaklines=true,
  frame=single,
  showstringspaces=false,
  tabsize=2,
}
\lstset{style=chemexcode}
\setlist{itemsep=2pt,topsep=4pt}
\setcounter{tocdepth}{2}

\title{Chemex 使用手册\\\large ChemDraw 与 LaTeX 的化学绘图桥接}
\author{Chemex 项目}
\date{版本 0.1.0；2026 年 11 月}

\begin{document}
\maketitle

\begin{abstract}
Chemex 将 ChemDraw 图稿中的动态化合物字段与 LaTeX 化合物编号连接起来，并为论文准备带编号的 CDXML、ChemDraw 矢量图和 MestReNova 谱图。原始 ChemDraw 文件仍由 ChemDraw 编辑；LaTeX 负责编号和论文排版；Lattice 可作为图稿浏览、项目管理与构建宿主。本手册说明当前 Chemex 的字段语法、编号规则、LaTeX 命令、命令行工作流、外部工具及适用边界。
\end{abstract}

\tableofcontents
\newpage

\section{简介}

\subsection{Chemex 做什么}

Chemex 在 CDXML 普通文字对象中识别动态字段，例如 \verb|[[@compound-001]]|。构建时，字段会替换为 LaTeX 计数器生成的纯文本编号。

LaTeX 正文用 \verb|\chemex{compound-001}| 引用同一化合物。改变计数器格式或插入图稿后，正文引用和 ChemDraw 图稿中的编号可由同一编号表更新。

Chemex 不重绘分子，也不把部分解析后的模型写回原始图稿。字段替换、原生导出和缓存均使用生成副本；源 CDXML/CDX 不会因编号或导出被改写。

\subsection{与 chemnum 的关系}

Chemex 借鉴化合物命名宏包手册常见的组织方式，并提供列表、子化合物和按需引用等相似概念；它不是 \texttt{chemnum} 的兼容实现。

\texttt{chemnum} 常见子标签写法是 \verb|\cmpd{main.sub}|，列表用逗号分隔。Chemex 的子字段写作 \verb!parent|child!；多个正文 key 可写在 \verb|\chemex{key1, key2}| 中。

Chemex 的关键区别是编号来源：默认计数器按论文中的图稿字段顺序推进，而不是简单按正文中首次提到某个 key 的位置编号。父子关系在 CDXML 文本对象内首次出现时建立，无须在导言区登记每个子项。计数器本身仍按 LaTeX 习惯在导言区定义。

\section{安装与准备}

\subsection{当前发行形态}

Chemex 目前作为 Lattice 仓库中的 CLI、LaTeX 宏包和共享核心提供；不是独立的 CTAN 安装包。独立命令由仓库的 \texttt{chemex-cli} crate 构建。运行论文构建时，Chemex 会为 LaTeX 提供对应版本的 \texttt{chemex.sty} 和准备文件，不需要用户把内部辅助宏复制进论文。

在仓库根目录构建：

\begin{lstlisting}
cargo build --manifest-path src-tauri/chemex/Cargo.toml -p chemex-cli --locked
\end{lstlisting}

也可按下文使用 \texttt{cargo run} 直接执行命令。若通过 Lattice 使用 Chemex，通常从论文项目的 Build、Chemistry 或绘图预览界面操作，无须手动运行 CLI。

\subsection{工具需求}

\begin{itemize}
  \item 仅检查 CDXML：Chemex CLI 即可；不需要启动 ChemDraw。
  \item 生成带编号 CDXML：需要可用的 \texttt{latexmk} 和论文所需的 LaTeX 发行版。
  \item 导出 ChemDraw 矢量图：需要 Windows、已安装并注册自动化接口的 ChemDraw，以及 Inkscape。Chemex 调用 ChemDraw 导出 SVG，再由 Inkscape 转为矢量 PDF。
  \item 处理 Mnova 文件：需要 MestReNova。Windows 下可设置 \texttt{CHEMEX\_MESTRENOVA} 指向其可执行文件；未设置时按系统可用信息查找。
  \item 运行 \texttt{structures}：需要 \texttt{uv}；首次执行会在隔离环境中取得 Python 3.13 和固定版本 RDKit 2026.9.2。可用 \texttt{CHEMEX\_UV} 指定 \texttt{uv} 路径。
\end{itemize}

外部程序可加入 \texttt{PATH}，或通过以下变量指定：

\begin{center}
\begin{tabular}{ll}
\toprule
变量 & 用途 \\
\midrule
\texttt{CHEMEX\_LATEXMK} & \texttt{latexmk} 可执行文件 \\
\texttt{CHEMEX\_INKSCAPE} & Inkscape 可执行文件 \\
\texttt{CHEMEX\_MESTRENOVA} & MestReNova 可执行文件 \\
\texttt{CHEMEX\_UV} & \texttt{uv} 可执行文件 \\
\bottomrule
\end{tabular}
\end{center}

\section{快速开始}

\subsection{在 ChemDraw 中建立字段}

在 ChemDraw 中创建普通文字对象，并输入字段占位符：

\begin{lstlisting}
[[@product-001]]
\end{lstlisting}

将字段放在对应结构附近，保存为 CDXML。若希望 \texttt{structures} 将字段与结构关联，应把结构和其独立编号文字放在同一个 ChemDraw Group 中；关联规则见第~\ref{sec:groups}~节。

\subsection{编写 LaTeX 文稿}

\begin{lstlisting}
\documentclass{article}
\usepackage{chemex}
\newcounter{compound}
\renewcommand{\thecompound}{\arabic{compound}}

\begin{document}
Product \chemex{product-001} is shown below.
\includegraphics[width=0.8\linewidth]{figures/reaction.cdxml}
\end{document}
\end{lstlisting}

将字段 key 与 LaTeX 引用中的 key 完全一致。对于 Chemex 默认计数器，图稿字段按文稿中图稿的出现顺序编号；正文中的 \verb|\chemex| 引用读取相应编号，不会因为引用提早出现在图稿之前而抢占默认编号。

\subsection{运行 CLI}

将 \texttt{MANUSCRIPT.tex} 替换为论文路径；目标输出目录必须尚不存在，且其父目录必须已经存在。

\begin{lstlisting}
chemex render MANUSCRIPT.tex output/rendered
\end{lstlisting}

\texttt{render} 生成带编号的 CDXML 副本。需要同时导出图稿 PDF，使用 \texttt{export}；需要生成最终论文 PDF，使用 \texttt{build}：

\begin{lstlisting}
chemex export MANUSCRIPT.tex output/exported
chemex build  MANUSCRIPT.tex output/paper
\end{lstlisting}

\section{LaTeX 命令}

\subsection{化合物引用 \texttt{\textbackslash chemex}}

基本语法如下：

\begin{lstlisting}
\chemex[<counter>]{<key-list>}
\end{lstlisting}

方括号可省略；省略时使用 \texttt{compound} 计数器。

\begin{center}
\begin{tabularx}{\linewidth}{lXl}
\toprule
写法 & 作用 & 要求 \\
\midrule
\texttt{\textbackslash chemex\{product-001\}} & 引用一个化合物 & key 已由对应流程编号 \\
\texttt{\textbackslash chemex[scheme]\{product-001\}} & 使用自定义计数器 & 计数器须先定义 \\
\texttt{\textbackslash chemex\{a, b, c\}} & 引用多个化合物 & key 之间用逗号分隔 \\
\bottomrule
\end{tabularx}
\end{center}

列表按输入顺序输出，不进行编号排序、连续范围压缩或重复项去重。key 不能含逗号；每个 key 由 Chemex 单独查找。默认两个条目以英语 \texttt{and} 连接，三个及以上以逗号分隔并在末项前加 \texttt{and}。可在导言区调整列表标点和连接词：

\begin{lstlisting}
\renewcommand{\ChemexListSeparator}{、}
\renewcommand{\ChemexListPairSeparator}{和}
\renewcommand{\ChemexListFinalSeparator}{和}
\end{lstlisting}

列表的排版分隔符设置在导言区；父子化合物的关系不需要在导言区预先声明。

\subsection{ChemDraw 图稿 \texttt{\textbackslash chemd}}

\begin{lstlisting}
\chemd{figures/reaction.cdxml}
\chemd[compound]{figures/reaction.cdxml}
\includegraphics[width=0.8\linewidth]{figures/reaction.cdxml}
\end{lstlisting}

\verb|\chemd{file.cdxml}| 使用默认 \texttt{compound} 计数器；\verb|\chemd[counter]{file.cdxml}| 指定 LaTeX counter。\texttt{\textbackslash includegraphics} 也可直接引用 CDXML，并保留 \texttt{graphicx} 的常用图片选项，例如 \texttt{width}、\texttt{height} 和 \texttt{trim}。Chemex 会为编号和渲染准备私有副本，LaTeX 源文件继续引用原 CDXML 路径。

\subsection{谱图 \texttt{\textbackslash chemn} 与 \texttt{\textbackslash chemd}}

\verb|\chemn| 专用于 Mnova 文件；\verb|\chemd| 根据文件扩展名分派 CDXML 或 Mnova。默认选择 Mnova 文件中的第 1 张谱图。

\begin{lstlisting}
\chemn{figures/sample.mnova}
\chemd{figures/sample.mnova}
\chemn[2]{figures/sample.mnova}
\chemd[2]{figures/sample.mnova}
\chemn[spec=2]{figures/sample.mnova}
\end{lstlisting}

数字方括号形式 \verb|[2]| 是选择第 2 张谱图的旧式简写。推荐使用键值参数；可选参数如下：

\begin{center}
\begin{tabularx}{\linewidth}{lXl}
\toprule
参数 & 含义 & 默认值 \\
\midrule
\texttt{spec} & 谱图序号，从 1 开始 & \texttt{1} \\
\texttt{cmpd} & ChemDraw 字段中的化合物 key，用于查找分子式 & 空 \\
\texttt{x} & 分子式左上角相对谱图宽度的位置百分比 & \texttt{50} \\
\texttt{y} & 分子式左上角相对谱图高度的位置百分比 & \texttt{15} \\
\texttt{size} & 分子式字号；支持 \texttt{pt}、\texttt{bp}、\texttt{mm}、\texttt{cm}、\texttt{in} & \texttt{8pt} \\
\bottomrule
\end{tabularx}
\end{center}

\texttt{x}、\texttt{y} 是 0--100 的整数百分比，坐标原点位于谱图左上角。分子式会缩放并限制在谱图图片范围内；叠加层不改变谱图在 LaTeX 中的图片盒尺寸。

\begin{lstlisting}
\chemn[spec=2,cmpd=product-17,x=78,y=12,size=8pt]{figures/sample.mnova}
\end{lstlisting}

\texttt{cmpd} 必须指向可关联到分子结构的有效 key；公式不可用或无法解析时，构建会报告错误，而不会猜测公式。

\section{编号规则}

\subsection{定义计数器与显示格式}

\texttt{chemex.sty} 使用默认计数器 \texttt{compound}，但不会替用户创建该计数器。应在导言区先定义计数器，再设置 \verb|\thecompound| 的格式：

\begin{lstlisting}
\newcounter{compound}
\renewcommand{\thecompound}{\arabic{compound}}
\end{lstlisting}

可使用 LaTeX 原生格式命令，例如 \verb|\arabic|、\verb|\roman|、\verb|\Roman|、\verb|\alph| 和 \verb|\Alph|。任意合法的 \texttt{\textbackslash the<counter>} 重定义也可使用，只要最终展开为 ChemDraw 可显示的普通文本。

\subsection{按章节编号}

在 \texttt{report} 或 \texttt{book} 文档中，将计数器设置为随章节重置，并把章节号放入显示格式：

\begin{lstlisting}
\documentclass{report}
\usepackage{chemex}
\newcounter{compound}[chapter]
\renewcommand{\thecompound}{\thechapter-\arabic{compound}}
\end{lstlisting}

第 2 章第 1 个父化合物及其子项会显示为 \texttt{2-1}、\texttt{2-1a}、\texttt{2-1b}。子字母附加在完整父编号之后，因此章节前缀不会与子编号冲突。\texttt{article} 文档没有 \texttt{chapter} counter，可按节或其他自定义计数器组织编号。

\subsection{默认与自定义 counter 的时序}

默认 \texttt{compound} 编号以图稿字段的原文顺序为基准。第一次遇到某个父 key 时推进计数器；相同 key 再次出现复用原号。正文中的 \verb|\chemex{key}| 可出现在图稿之前，但不会让默认编号抢先分配。

使用自定义计数器，如 \verb|\chemd[scheme]{...}| 或 \verb|\chemex[scheme]{key}|，该计数器按 LaTeX 执行顺序工作；首次正文引用或图稿字段出现时推进。counter 通过 \verb|\newcounter| 声明，重置使用 \verb|\setcounter|。某个 key 一旦绑定到一个 counter，就不能再以另一个 counter 使用。

字段 key 是化合物身份，不是每次出现时临时生成的编号。重置 counter 不会清除同一 key 已绑定的编号；编号给新化合物时，应在 ChemDraw 中为它使用新 key。不同年份、章节或不同 counter 管理的化合物也应使用不同 key。

\section{ChemDraw 动态字段}

\subsection{普通字段语法}

在 ChemDraw 的普通文字对象中输入占位符。生成图稿时，Chemex 将方括号和 key 替换为 LaTeX 计算出的编号：

\begin{lstlisting}
[[@compound-001]]
\end{lstlisting}

这是 Chemex 的构建期文本字段，不是 ChemDraw 自带的自动编号域；在 ChemDraw 编辑源文件时，标记保持原样。编号后的结果只写入 Chemex 生成的 CDXML 副本。

普通 key 可含 ASCII 字母、数字、下划线、点、冒号和连字符。逗号用于 LaTeX 列表分隔，不能出现在 key 内。保留简单、稳定且便于检索的 key，例如 \texttt{substrate-01}、\texttt{product-17} 或 \texttt{paperA.chapter2.03}。

\subsection{带前后缀的字段}

若要把化合物编号嵌入较长文字，可将待替换 key 部分放在下划线之后的前缀和插入符号之前的后缀之间：

\begin{lstlisting}
[[@product-001_ Scheme ^ product ]]
\end{lstlisting}

编号替换后得到类似 \texttt{Scheme 3 product} 的文字。前缀和后缀可以为空，也可含空格；替换仅覆盖 key 所在位置，装饰和样式保留。使用这种装饰形式时，key 本身不能包含下划线。

\subsection{父化合物与子化合物}

Chemex 用一个竖线分隔父 key 和子 key：

\begin{lstlisting}
[[@compound-7]]
[[@compound-7|isomer-a]]
[[@compound-7|isomer-b]]
\end{lstlisting}

Chemex 无须在导言区声明父子关系。构建时按图稿字段首次出现的顺序为每个父项的子 key 分配 \texttt{a}、\texttt{b}、\ldots{}；重复的同一父子组合沿用原字母。父编号推进一次，子项不推进主 counter。

只需要在图中显示子编号时，可以不放单独的父字段：首个 \verb!parent|child! 会自动建立父编号。若还需显示父编号，再添加 \verb|[[@parent]]|。正文引用子项使用同一复合 key：

\begin{lstlisting}
\chemex{compound-7|isomer-a} and \chemex{compound-7|isomer-b}
\chemex{compound-7|isomer-a, compound-7|isomer-b}
\end{lstlisting}

完整父子 key 的唯一身份是 \texttt{parent|child} 组合；不同父项可各自使用相同的子 key。每个父 key 最多有 26 个不同子 key；当前只支持父项加一层子项。

\subsection{字段识别范围与排版}

Chemex 仅替换可识别的普通 CDXML 文字对象。原子标签、未知对象文字、属性、注释和 \texttt{ObjectTag} 元数据不是动态字段。字段可跨 ChemDraw 的多个文字样式片段；编号采用 key 起始处的样式，前后文字和其他样式片段保留。

单独占满一个文字对象的字段被视为独立编号，生成副本会尝试以原文字框的水平中心居中，并保留基线。没有可靠位置/文字框或存在旋转时，Chemex 会给出对齐诊断并保留原排版，但仍可替换编号。句子中的内嵌字段不强制居中。

\subsection{结构分组与分析}
\label{sec:groups}

若需要让 \texttt{structures} 把编号与结构联系起来，在 ChemDraw 中选中结构和对应的独立编号文字，使用 ChemDraw 的 Group 命令组成同一组。组内可有多个结构片段，例如盐的组分；Chemex 保留这些片段，不自动拆盐或判断化学等价性。

一个组只能对应一个不同的化合物 key。同 key 的重复标签可以共存；不同 key 放在同组会报告歧义，Chemex 不猜测哪个字段代表结构。未分组字段仍可编号和替换，只是没有结构绑定。

\section{命令行参考}

\subsection{命令一览}

\begin{center}
\begin{tabularx}{\linewidth}{lX}
\toprule
命令 & 用途 \\
\midrule
\texttt{chemex inspect SOURCE.cdxml} & 检查 CDXML 对象、字段、分组与解析诊断，输出 JSON \\
\texttt{chemex structures SOURCE.cdxml} & 分析已关联结构，输出分子信息与重复项报告 \\
\texttt{chemex render MANUSCRIPT.tex NEW\_DIR} & 运行 TeX 编号流程并生成编号 CDXML 副本 \\
\texttt{chemex export MANUSCRIPT.tex NEW\_DIR} & 准备编号图稿并导出 ChemDraw/Mnova PDF \\
\texttt{chemex build MANUSCRIPT.tex NEW\_DIR} & 构建论文，最终目录只发布论文 PDF \\
\bottomrule
\end{tabularx}
\end{center}

不带参数或使用 \texttt{--help} 可查看命令摘要。当前命令不提供通用的覆盖输出选项；目标目录已存在时，命令会失败并要求换一个新目录。

\subsection{检查 CDXML：\texttt{inspect}}

\begin{lstlisting}
chemex inspect figures/reaction.cdxml
\end{lstlisting}

输出 JSON 报告，包括对象清单、字段 key、文字对象索引、分组/片段关联以及解析诊断。对象索引只在当前文件中有效，不是永久 ChemDraw ID，也不应被当作化合物身份。

\subsection{结构分析：\texttt{structures}}

\begin{lstlisting}
chemex structures figures/reaction.cdxml
\end{lstlisting}

该命令只分析已由字段和 ChemDraw Group 关联的结构。RDKit 解析后报告分子式、规范异构 SMILES 和在条件满足时的精确重复项。它不是完整 CDXML 渲染器，也不做盐拆分、质子化态归一、互变异构或未指定立体化学的猜测。查询结构、解析/映射不确定或分组歧义的条目不会自动判重。

\subsection{编号副本：\texttt{render}}

\begin{lstlisting}
chemex render manuscript.tex .chem/render-01
\end{lstlisting}

Chemex 通过 TeX 读取论文 counter 和编号格式，输出替换字段后的 CDXML 及编号过程文件。此命令不发布最终论文 PDF。目标目录必须不存在；源 TeX 和 CDXML 不会被覆盖。

\subsection{图稿导出：\texttt{export}}

\begin{lstlisting}
chemex export manuscript.tex .chem/export-01
\end{lstlisting}

除编号准备外，此命令还为 ChemDraw 图稿生成 SVG/PDF，为选定的 Mnova 谱图生成 PDF。ChemDraw 当前通过本机原生 SVG 导出，再由 Inkscape 转为矢量 PDF；不会用位图重绘结构。失败时保留诊断日志和必要中间产物，不把不完整结果作为成功输出发布。

若希望之后自行运行 \texttt{latexmk}，可将导出目录放到 \texttt{TEXINPUTS} 搜索路径最前面。目录中包含准备好的 \texttt{chemex.sty}、注册表和图稿 PDF；LaTeX 源文件仍引用原来的 CDXML/Mnova 路径：

\begin{lstlisting}
$export = (Resolve-Path .\output\exported).Path
$env:TEXINPUTS = "$export;$env:TEXINPUTS;"
latexmk -pdf -interaction=nonstopmode manuscript.tex
\end{lstlisting}

末尾分号保留 TeX 发行版默认搜索路径。重新导出时使用一个新的目标目录，随后更新 \texttt{TEXINPUTS} 指向该目录。

\subsection{完整论文：\texttt{build}}

\begin{lstlisting}
chemex build manuscript.tex .chem/paper-01
\end{lstlisting}

执行 Chemex 准备、图稿导出和 LaTeX 构建。成功目标目录只包含最终论文 PDF；中间文件由独立 CLI 管理。构建失败返回非零退出码，并保留诊断信息。CLI 的完整项目依赖发现范围有限，复杂多文件项目应使用 Lattice 原生 Build。

\section{输入、缓存与安全}

\subsection{输入和输出规则}

\begin{itemize}
  \item CDXML 输入须为 UTF-8，可带 BOM；CDX 二进制文件应先用 ChemDraw 另存为 CDXML。
  \item 单个 CDXML 源文件最多 64 MiB，XML 节点最多 250{,}000 个。
  \item CLI 会间隔 100 毫秒读取源文件两次并比较内容；检测到写入中的变化时拒绝继续。这是稳定读取检查，不是文件锁或持续监控。
  \item 输出目标目录不得已存在，父目录须预先创建。重新运行时指定新的目标目录。
  \item 字段解析与替换可以处理多页 CDXML；ChemDraw 原生图像导出目前拒绝未验证的多页/跨页画布布局。
\end{itemize}

\subsection{缓存}

独立 CLI 在主 \texttt{.tex} 所在项目的 \texttt{.chem/cache/} 下保存可复用导出：ChemDraw 位于 \texttt{chemdraw/}，Mnova 位于 \texttt{mestrenova/}。缓存依据源文件、编号后图稿及渲染器信息判断是否仍有效；字段编号改变会改变生成 CDXML 和相应渲染内容。Lattice 的化学分析数据库位于项目 \texttt{.chem/molecules.sqlite}。

缓存是派生产物，不是源稿。删除缓存会导致重算/重导；不要编辑缓存文件来修改权威图稿。应在 ChemDraw 源文件或 LaTeX 源文稿中修改内容。

\subsection{原件保护}

字段替换只对新字符串做局部补丁，并检查输出 XML 仍可解析。ChemDraw 导出脚本打开的是私有输入副本，只关闭自己打开的临时文档，不关闭用户已经打开的 ChemDraw 文档。外部阶段超时为 120 秒；超时时原生应用可能仍持有临时副本，应先查看日志和中间目录，再决定是否清理。

\section{常见问题}

\begin{description}[style=nextline]
  \item[正文显示 \texttt{??}] 检查 key 拼写是否与 CDXML 字段完全一致、编号字段是否存在，以及当前构建是否运行 Chemex 准备流程。默认 counter 下，正文引用不会替代缺失的图稿字段自动编号。
  \item[未知 LaTeX counter] 自定义编号通过 \verb|\chemex[counter]{key}| 或 \verb|\chemd[counter]{file.cdxml}| 指定。使用前先以 \verb|\newcounter{counter}| 定义 counter，并检查名称拼写。
  \item[同一 key 使用多个 counter] 同一个化合物身份不能跨 counter 复用。为新计数体系中的化合物选择新 key。
  \item[子编号错误或重复] 检查字段是否只含一个竖线、父 key 是否一致、子 key 是否误写，以及同父项是否超过 26 个不同子 key。重复出现相同的 \verb!parent|child! 是同一个子项，不会分配新字母。
  \item[字段未被替换] 确认它位于普通文字对象，而非原子标签、ObjectTag、XML 属性或未知对象中；检查字段括号完整、CDXML 语法有效。
  \item[结构没有关联] 在 ChemDraw 中将结构和编号文字分组；确认一个组内不存在多个不同 key。未关联不妨碍普通文本编号替换。
  \item[ChemDraw 导出失败] 检查 ChemDraw COM 自动化是否可用、是否支持 SVG 导出，以及 Inkscape 路径是否正确。ChemDraw 多页导出目前不支持。
  \item[谱图序号错误] Mnova 序号从 1 开始；确保文件含有所选谱图。检查 \texttt{CHEMEX\_MESTRENOVA} 或可执行文件发现路径。
  \item[目标目录已存在] CLI 不覆盖已有结果。改用新的输出目录；不要删除或移动不确定来源的目录。
\end{description}

\section{限制与兼容性}

当前应注意以下范围：

\begin{itemize}
  \item 父子化合物仅支持一层；每个父 key 最多 26 个子项。
  \item ChemDraw 多页/跨页画布的矢量导出尚未验证；字段替换支持多页，但这不等于导出兼容性已验证。
  \item Chemex 不替代 ChemDraw 的化学编辑或完整 CDXML 实现；RDKit 只负责有限的分子/反应分析。
  \item \texttt{ObjectTag} 可读取，但 ChemDraw 编辑保存往返后自定义身份信息是否保留，需按目标 ChemDraw 版本实测。
  \item \texttt{structures} 的精确重复判断要求结构可解析且立体化学信息明确；它不推断未知化学语义。
  \item 独立 CLI 的 \texttt{build} 不接管复杂的多文件依赖、BibTeX 依赖发现或任意自定义构建脚本；Lattice 原生 Build 用于完整宿主工作流。
\end{itemize}

\appendix
\section{语法速查}

\begin{lstlisting}
ChemDraw field:      [[@key]]
Decorated field:     [[@key_prefix^suffix]]
Child field:        [[@parent|child]]
Single reference:    \chemex{key}
Custom counter:      \chemex[counter]{key}
Compound list:       \chemex{key1, key2, key3}
CDXML drawing:       \chemd[counter]{drawing.cdxml}
Mnova spectrum:      \chemn[spec=2]{sample.mnova}
Spectrum + formula:  \chemn[spec=2,cmpd=key,x=78,y=12,size=8pt]{sample.mnova}
\end{lstlisting}

\section{命令速查}

\begin{lstlisting}
chemex --help
chemex inspect SOURCE.cdxml
chemex structures SOURCE.cdxml
chemex render MANUSCRIPT.tex NEW_DIRECTORY
chemex export MANUSCRIPT.tex NEW_DIRECTORY
chemex build MANUSCRIPT.tex NEW_DIRECTORY
\end{lstlisting}

\section{项目与参考}

Chemex 源码和独立示例位于 Lattice 仓库的 \texttt{src-tauri/chemex/}。项目 README 提供快速示例；\texttt{docs/DEVELOPMENT.md} 记录实现约束、字段契约和验证边界；\texttt{examples/} 提供 LaTeX 示例。

本手册的章节组织参考常见 LaTeX 宏包用户手册；\texttt{chemnum} 官方项目及其手册可见：\url{https://github.com/cgnieder/chemnum}。Chemex 与其命令名、分隔符和编号时序并不兼容。

\end{document}
